%% dssyk_krylov_FRW_extension.tex
%% COMPANION NOTE: Krylov Complexity in the FRW Diamond
%% Extending Heller-Papalini-Schuhmann (arXiv:2412.17785) to FRW
%%
%% Status: DRAFT 2026-05-04
%% Authors: (to be determined)
%%
%% Anti-hallu notes (embedded throughout):
%%   - HPS 2412.17785 is about DSSYK/sine-dilaton gravity (C=V matching),
%%     NOT algebraic-KMS Krylov theory. Its results are cited accurately below.
%%   - Parker et al. 1812.08657: abstract gives lambda_L <= 2alpha, not
%%     explicitly lambda_K = 2pi/beta; SYK saturation is numerical.
%%   - FRW type-II_inf is established in frw_note.tex (Thm 1/2, 2026-05-02).
%%   - ECI v6.0.44 (sec:mss-shadow) flags the bridge S_gen <-> C_Krylov as
%%     a working conjecture, not a theorem.
%%   - All [CONJECTURE] tags are mandatory.

\documentclass[12pt,a4paper]{article}

\usepackage{amsmath,amssymb,amsthm}
\usepackage{hyperref}
\usepackage{cite}
\usepackage{geometry}
\usepackage{xcolor}
\geometry{margin=2.5cm}

%% Theorem environments
\newtheorem{theorem}{Theorem}
\newtheorem{proposition}[theorem]{Proposition}
\newtheorem{lemma}[theorem]{Lemma}
\newtheorem{corollary}[theorem]{Corollary}
\newtheorem{conjecture}[theorem]{Conjecture}
\theoremstyle{definition}
\newtheorem{remark}[theorem]{Remark}
\newtheorem{question}[theorem]{Open Question}

%% Macros
\newcommand{\CK}{\mathcal{C}_K}
\newcommand{\AO}{\mathcal{A}(\mathcal{O})}
\newcommand{\lK}{\lambda_K}
\newcommand{\lL}{\lambda_L}
\newcommand{\bw}{\mathrm{BW}}
\newcommand{\HL}{\mathrm{HL}}
\newcommand{\FRW}{\mathrm{FRW}}
\newcommand{\Mink}{\mathrm{M}}
\newcommand{\sgen}{S_{\mathrm{gen}}}
\newcommand{\Smod}{\mathcal{C}_K^{\mathrm{mod}}}
\newcommand{\KMS}{\mathrm{KMS}}
\newcommand{\dS}{\mathrm{dS}}
\newcommand{\TODO}[1]{\textcolor{red}{[TODO: #1]}}

\title{\textbf{Krylov Complexity in the FRW Diamond}\\[4pt]
       \large Extending Heller--Papalini--Schuhmann to Non-Stationary Cosmology
       via Conformal Pullback}

\author{\textit{(Author names to be determined)}\\[2pt]
  \small \textit{Companion note to: ``The Modular Shadow Conjecture''
  [Letters in Mathematical Physics, submitted]}\\
  \small \textit{and to: ``Type II$_\infty$ for the gravitational algebra in
  radiation-dominated FRW'' \cite{FRWnote}}}

\date{Preprint: 2026-05-04}

\begin{document}

\maketitle

\begin{abstract}
We extend the analysis of Heller, Papalini, and Schuhmann~\cite{HPS2024}
--- which establishes a quantitative match between Krylov spread complexity
in double-scaled SYK (DSSYK) and the complexity-equals-volume proposal in
sine-dilaton gravity at finite temperatures --- to the FRW causal diamond.
The algebraic foundation is provided by the type-II$_\infty$ classification
of the FRW diamond algebra~\cite{FRWnote}: for the conformally coupled
massless scalar in radiation- or matter-dominated FRW, the crossed-product
algebra $\mathcal{A}(\mathcal{O})_\FRW \rtimes_\sigma \mathbb{R}$ is type
II$_\infty$, with the modular automorphism group unitarily intertwined with
the Hislop--Longo flow on the conformally identified Minkowski diamond.
The KMS state at $\beta = 2\pi$ (conformal pullback of Bisognano--Wichmann)
gives a modular Hamiltonian with continuous unbounded spectrum, enabling
a Lanczos recursion.
We conjecture [CONJECTURE] that the Krylov complexity growth rate in the
FRW diamond satisfies $\lK^\FRW = 2\pi/\beta = 1$ in natural units,
saturating the Maldacena--Shenker--Stanford (MSS) chaos bound, as a
structural consequence of the type-II$_\infty$ classification and the
KMS condition.
The central bridge $\sgen \leftrightarrow \Smod$ is flagged explicitly
as unproven; the note records supporting arguments, identifies the obstructions,
and proposes a programme of $\sim 4$ weeks of computation to either
substantiate or refute the conjecture.
\end{abstract}

\tableofcontents

%% ─────────────────────────────────────────────────────────────────────
\section{Background and Motivation}
\label{sec:background}
%% ─────────────────────────────────────────────────────────────────────

\subsection{Krylov complexity and the MSS chaos bound}

Parker, Cao, Avdoshkin, Scaffidi, and Altman~\cite{Parker2019}
introduced \emph{Krylov} (Lanczos) complexity $\CK(t)$ as a
basis-independent measure of operator growth.
Given an operator $\mathcal{O}$ evolving under a Hamiltonian $H$,
one constructs iteratively the orthonormal \emph{Krylov basis}
$\{|O_n\rangle\}$ from the nested commutators $\mathcal{L}^n\mathcal{O}$
via the Gram--Schmidt procedure, where $\mathcal{L} = [H, \cdot]$ is the
Liouvillian superoperator.
The Krylov complexity is the mean position in this basis:
$\CK(t) = \sum_n n \, |\langle O_n | \mathcal{O}(t)\rangle|^2$.
The key quantity in the recursion is the sequence of \emph{Lanczos
coefficients} $\{b_n\}_{n \ge 1}$, defined by the three-term recurrence.

\textbf{Universal operator growth hypothesis~\cite{Parker2019}.}
\emph{In generic (non-integrable) many-body systems, the Lanczos
coefficients $b_n$ grow asymptotically linearly:}
\begin{equation}
  b_n \;\sim\; \alpha\, n \;+\; \text{(sub-leading)},
  \qquad n \to \infty,
  \label{eq:lanczos-linear}
\end{equation}
\emph{where $\alpha > 0$ is a model-dependent rate, with a logarithmic
correction in 1+1 dimensions.
The rate $\alpha$ upper-bounds the Lyapunov exponent:
$\lL \le 2\alpha$.
Since the universal low-temperature bound gives $\lL \le 2\pi T = 2\pi/\beta$
(MSS), saturation of both bounds requires $\alpha = \pi/\beta$ and
$b_n \sim (\pi/\beta)\,n$ in maximally chaotic systems.}

The MSS chaos bound~\cite{MSS2016} states that in any large-$N$ quantum
system with a thermal density matrix at inverse temperature $\beta$,
the OTOC Lyapunov exponent satisfies
\begin{equation}
  \lL \;\le\; \frac{2\pi}{\beta}.
  \label{eq:mss}
\end{equation}
Equality holds in maximally chaotic systems including near-extremal
black holes and the SYK model.
The modular Krylov complexity version of MSS saturation, established
numerically in SYK at large $N$~\cite{Parker2019}, is
$\lK = 2\pi/\beta$ at saturation.

\subsection{Heller--Papalini--Schuhmann: DSSYK and holographic Krylov}

The paper by Heller, Papalini, and Schuhmann~\cite{HPS2024}
(arXiv:2412.17785, PRL 135 151602, 2025) establishes a quantitative
holographic match:

\begin{quote}
\textbf{Verbatim abstract of~\cite{HPS2024}:}
``One of the important open problems in quantum black hole physics is
a dual interpretation of holographic complexity proposals.
To date the only quantitative match is the equality between the Krylov
spread complexity in triple-scaled SYK at infinite temperature and the
complexity = volume proposal in classical JT gravity.
Our work utilizes the recent connection between double-scaled SYK and
sine-dilaton gravity to show that the quantitative relation between Krylov
spread complexity and complexity = volume extends to finite temperatures
and to full quantum regime on the gravity side at disk level.
From the latter we isolate the first quantum correction to the
complexity = volume proposal and propose to view it as a complexity of
quantum fields in the bulk.
Finally, we comment on the switchback effect, whose presence would make
the Krylov spread complexity a fully fledged holographic complexity at
least in sine-dilaton gravity.''
\end{quote}

\begin{remark}[Scope of HPS result]
\label{rem:HPS-scope}
The HPS result~\cite{HPS2024} is a \emph{holographic} result about
DSSYK/sine-dilaton gravity, not a derivation from algebraic quantum field
theory or type-II crossed-product formalism.
It establishes:
(a) a quantitative $C = V$ match at finite temperatures in DSSYK;
(b) the first quantum correction to $C = V$ as bulk-field complexity;
(c) evidence for the switchback effect in DSSYK.
It does \emph{not} directly prove the Krylov exponent in a type-II$_\infty$
algebraic setting.
The extension to FRW proposed in this note therefore requires additional
algebraic input (the type-II$_\infty$ classification of the FRW diamond,
\S\ref{sec:FRW-setup}) and an explicit conjecture [CONJECTURE] connecting
HPS's holographic Krylov to the algebraic FRW modular flow.
\end{remark}

\subsection{Modular shadow and MSS saturation (ECI sec:mss-shadow)}

The ECI programme~\cite{ECI_v6} identifies the MSS chaos bound as a
``modular shadow'' of the saturation of the generalised entropy growth rate
in the type-II$_\infty$ algebra.
The ECI saturation identity is (eq.~(20) of~\cite{ECI_v6}):
\begin{equation}
  \frac{1}{\sgen} \frac{d\sgen}{d\tau}\bigg|_{\mathrm{sat}}
  \;=\; \frac{2\pi T_H}{\hbar}
  \;=\; \frac{\kappa_R}{\hbar},
  \label{eq:eci-mss}
\end{equation}
where $\kappa_R = 2\pi T_H$ is the Bisognano--Wichmann surface gravity.
Under the Faulkner--Speranza identification $\sgen \leftrightarrow
\Smod$~\cite{FaulknerSperanza2024},
equation~\eqref{eq:eci-mss} produces the MSS bound at equality.
\textbf{This identification is a working conjecture, not a theorem.}
(ECI v6.0.44, \S sec:mss-shadow, explicitly states: ``We flag this
identification as a working conjecture rather than a theorem:
the bridge $\sgen \leftrightarrow \Smod$ is consistent with the
Witten--CLPW--Caputa--Faulkner--Speranza programme but has not been
derived as a theorem in type-II$_\infty$ crossed-product algebras.'')

The companion paper~\cite{ModularShadow} develops this into the
Modular Shadow Conjecture.
The present note asks: does the same structural picture hold in the
\emph{FRW diamond}, where no global Killing vector exists?

%% ─────────────────────────────────────────────────────────────────────
\section{The Heller--Papalini--Schuhmann Setting and Its Scope}
\label{sec:HPS}
%% ─────────────────────────────────────────────────────────────────────

\subsection{DSSYK and sine-dilaton gravity}

Heller, Papalini, and Schuhmann~\cite{HPS2024} work in the setting of
\emph{double-scaled SYK} (DSSYK), a variant of the Sachdev--Ye--Kitaev
model defined by the limit $N, q \to \infty$ with $\lambda = 2 \log q / N$
fixed.
In this regime the model is dual, via the chord-diagram expansion, to
\emph{sine-dilaton gravity} on a two-dimensional background.
The Krylov spread complexity in DSSYK at inverse temperature $\beta$
is given by:
\begin{equation}
  2|\log q|\,\CK(t)_\beta \;=\; \langle \hat{L} \rangle
  \;=\; \bigl[-\partial_\Delta \langle e^{-\Delta\hat{L}}\rangle\bigr]_{\Delta=0},
  \label{eq:HPS-main}
\end{equation}
where $\hat{L}$ is the chord-length operator in sine-dilaton gravity.
The central result of~\cite{HPS2024} is that
$\CK(t)_\beta = \langle \hat L \rangle_\beta / (2|\log q|)$
matches the complexity-equals-volume (C=V) proposal in sine-dilaton gravity
at finite temperatures and at disk-level in the quantum $1/C_\phi$ expansion.

\subsection{What is AdS/static-specific vs. what can transport to FRW}

The HPS analysis relies on:
\begin{enumerate}
\item \textbf{Static patch / thermal state.}
  DSSYK at temperature $\beta^{-1}$ corresponds to a static thermal
  state on the JT/sine-dilaton gravity side. There is an effective
  time-translation symmetry (the Hamiltonian of the SYK model)
  generating the thermal KMS flow.

\item \textbf{Chord diagram / matrix model structure.}
  The quantitative $C = V$ match uses chord combinatorics specific
  to DSSYK. These are not available in a general curved background.

\item \textbf{Holographic dictionary (JT/DSSYK).}
  The gravity dual is two-dimensional and its complexity = volume
  proposal is specific to dilaton gravity theories.
\end{enumerate}

\textbf{What survives in FRW:}
The \emph{structural} content of the HPS result --- that Krylov
spread complexity is a holographically computable measure of operator
growth that saturates the MSS bound at $\lK = 2\pi/\beta$ ---
is what we wish to transport to the FRW diamond.
The transport is possible because:
\begin{itemize}
\item The FRW diamond algebra $\mathcal{A}(\mathcal{O})_\FRW$ is
  type-II$_\infty$ (established, \S\ref{sec:FRW-setup}),
  the same algebraic class as the AdS/JT algebra at large $N$.
\item The KMS state on the FRW diamond at $\beta = 2\pi$ is the
  conformal pullback of the Minkowski vacuum (Hislop--Longo), which
  is structurally analogous to the Hartle--Hawking state in AdS.
\item The Lanczos recursion is well-defined in any von Neumann algebra
  equipped with a faithful normal state (the modular Krylov construction
  of Caputa et al.~\cite{Caputa2024}).
\end{itemize}

\textbf{What does not transport:}
\begin{itemize}
\item The quantitative $C = V$ match depends on the chord-diagram
  structure and the holographic dictionary of DSSYK. No FRW analogue
  of this dictionary exists.
\item The quantum correction (bulk-field complexity) isolated by HPS
  is JT-gravity specific; in the FRW setting one would need a two-
  dimensional gravity dual for the FRW diamond, which is not available.
\item There is no FRW analogue of the global Hamiltonian $H_\mathrm{SYK}$
  generating the thermal flow; the modular Hamiltonian $K_\FRW$ (a
  non-local operator) must replace it.
\end{itemize}

%% ─────────────────────────────────────────────────────────────────────
\section{FRW Diamond Setup}
\label{sec:FRW-setup}
%% ─────────────────────────────────────────────────────────────────────

\subsection{The metric and the comoving observer}

We work in spatially flat ($k=0$) FRW spacetime in conformal time $\eta$.
The Hadamard property of the SLE state on the BKL attractor for Bianchi~IX
(a related non-stationary setting) is established in~\cite{BIX_lemma_A1};
the type-II$_\infty$ analogy is an instructive comparison.


\begin{equation}
  ds^2_\FRW \;=\; a(\eta)^2\bigl(-d\eta^2 + d\mathbf{x}^2\bigr),
  \qquad \mathbf{x}\in\mathbb{R}^3.
  \label{eq:FRW-metric}
\end{equation}
For radiation-dominated matter content ($p = \rho/3$), the Friedmann
equations give $a(\eta) = a_0\,\eta$; for matter-dominated content
($p = 0$), $a(\eta) = a_0\,\eta^2$.
Both cases are conformally flat, with the Minkowski metric
$\eta_{\mu\nu}=\mathrm{diag}(-1,+1,+1,+1)$ as conformal target.

A comoving observer at $\mathbf{x} = 0$ on $\eta \in (\eta_i, \eta_f)$
(with $0 < \eta_i < \eta_f < \infty$) sweeps out a \emph{causal diamond}
\begin{equation}
  D(\mathcal{O}) \;=\; \bigl\{(\eta,\mathbf{x})\;:\;
  |\mathbf{x}| < \eta_f - \eta,\;\eta > \eta_i + |\mathbf{x}|\bigr\}.
\end{equation}
In the conformally identified Minkowski coordinates, $D(\mathcal{O})$
is the same set of coordinate points as a doubly-bounded Minkowski diamond
$M_\mathcal{O}$.

\textbf{No global Killing vector.}
Neither radiation-dominated nor matter-dominated FRW admits a
global timelike Killing vector field along a comoving worldline.
The CLPW construction~\cite{CLPW2023, Witten2022} and its extensions
to Killing horizons~\cite{FaulknerSperanza2024} are not directly
applicable.

\subsection{Conformally coupled scalar and the type-II$_\infty$ theorem}

We consider the conformally coupled massless scalar field:
\begin{equation}
  S[\varphi] \;=\; -\tfrac{1}{2}\int d^4x\,\sqrt{-g}
  \Bigl[g^{\mu\nu}\partial_\mu\varphi\,\partial_\nu\varphi
        \;-\; \tfrac{1}{6}\,R\,\varphi^2\Bigr],
  \qquad \xi = \tfrac{1}{6}.
  \label{eq:action}
\end{equation}
Under the Weyl rescaling $\tilde\varphi := a(\eta)\,\varphi$, the
field equation $(\Box_g - R/6)\varphi = 0$ maps exactly to
$\Box_\Mink \tilde\varphi = 0$ (conformal-covariance identity,
sympy-verified in~\cite{FRWnote}).
The Bunch--Davies (conformal) vacuum $\Omega_\FRW$ maps to the
Minkowski vacuum $\Omega_\Mink$ under the unitary
$U: \mathcal{H}_\FRW \to \mathcal{H}_\Mink$ implementing
$\varphi(f) \mapsto \tilde\varphi(a^3 \cdot f)$.

The following theorem is established in~\cite{FRWnote}:

\begin{theorem}[{Type II$_\infty$ for the FRW diamond, \cite{FRWnote} Theorem 1}]
\label{thm:FRW-typeII}
Let $\varphi$ be the conformally coupled massless scalar field on
radiation-dominated FRW with $a(\eta) = a_0\eta$, in the conformal
vacuum $\Omega_\FRW$.
Let $D_\mathcal{O}$ be a doubly-bounded causal diamond for a comoving
observer with $0 < \eta_i < \eta_f < \infty$.
Let $\sigma^\FRW$ denote the modular automorphism group of
$(\mathcal{A}(D_\mathcal{O})_\FRW, \Omega_\FRW)$.
Then:
\begin{enumerate}
\item[(1)] $\sigma_t^\FRW = \mathrm{Ad}(U^{-1}) \circ \sigma_t^\HL \circ \mathrm{Ad}(U)$,
  where $\sigma^\HL$ is the Hislop--Longo modular automorphism group of
  the conformally identified Minkowski diamond.
\item[(2)] The Connes spectrum satisfies $\Gamma(\sigma^\FRW) = \mathbb{R}$.
\item[(3)] The crossed product $\mathcal{A}(D_\mathcal{O})_\FRW
  \rtimes_{\sigma^\FRW} \mathbb{R}$ is a type II$_\infty$ factor.
\end{enumerate}
The same conclusion holds for matter-dominated FRW ($a(\eta) = a_0\eta^2$)
by the same proof.
\end{theorem}

\begin{remark}[No Killing vector required]
The proof of Theorem~\ref{thm:FRW-typeII} uses only:
(a) conformal flatness of $g_\FRW$;
(b) conformal invariance of the field equation ($\xi = 1/6$, $m = 0$);
(c) the Hislop--Longo theorem on the conformally identified Minkowski diamond;
(d) Tomita--Takesaki uniqueness (modular operators intertwine under
the unitary $U$);
(e) the Connes--Takesaki classification (Connes spectrum $\mathbb{R}$
implies type II$_\infty$ of the crossed product).
No timelike Killing vector of $g_\FRW$ is invoked.
This is the first example of the CLPW type-II$_\infty$ classification
in a genuinely non-stationary FRW background with no Killing symmetry.
\end{remark}

\subsection{The FRW modular Hamiltonian and its KMS state}

By Theorem~\ref{thm:FRW-typeII}(1), the FRW modular Hamiltonian is
\begin{equation}
  K_\FRW \;=\; U^{-1}\,K_\HL\,U,
  \label{eq:K-FRW}
\end{equation}
where $K_\HL = -\log\Delta_\HL$ is the Hislop--Longo modular Hamiltonian
for the Minkowski diamond.
The Hislop--Longo modular flow is the one-parameter M\"{o}bius group
of conformal transformations of the Minkowski diamond fixing its tips;
its generator $K_\HL$ is conjugate, via the conformal group
$\mathrm{SO}(2,4)$, to a boost generator and has continuous spectrum
filling $\mathbb{R}$.
Consequently:
\begin{itemize}
\item $K_\FRW$ has spectrum $\mathrm{Spec}(K_\FRW) = \mathbb{R}$
  (inherited from $K_\HL$ by unitary conjugation).
\item The modular automorphism group $\sigma_s^\FRW = e^{isK_\FRW}
  (\cdot) e^{-isK_\FRW}$ implements a flow on $\mathcal{A}(D_\mathcal{O})_\FRW$
  that corresponds, in the Minkowski picture, to the M\"{o}bius flow
  of the diamond, dressed by the conformal factor $a(\eta(s))/a(\eta)$
  (explicit formula in~\cite{FRWconformal}).
\item The state $\Omega_\FRW$ is a KMS state at inverse temperature
  $\beta = 2\pi$ with respect to the modular flow $\sigma_s^\FRW$,
  by Tomita--Takesaki theory.
\end{itemize}

The conformal weight factor $a(\eta(s))/a(\eta)$ is the key difference
from the static/AdS case.
On the diamond orbit $(\eta(s), x^1(s)) = (\eta\cosh(2\pi s) + x^1\sinh(2\pi s),
\; \eta\sinh(2\pi s) + x^1\cosh(2\pi s))$, this weight is:
\begin{equation}
  w_\FRW(s;\eta,x^1) \;:=\; \frac{a(\eta(s))}{a(\eta)}
  \;=\; \frac{\eta\cosh(2\pi s) + x^1\sinh(2\pi s)}{\eta}
  \quad (a = a_0\eta),
  \label{eq:conf-weight}
\end{equation}
which satisfies $w_\FRW(0;\eta,x^1) = 1$ and $w_\FRW(s) \to e^{2\pi s}$
as $s \to +\infty$ (on the wedge limit $x^1 \ll \eta$), recovering
the Bisognano--Wichmann exponential at large boost parameter.

\subsection{Scope and honest limitations}

The following limitations of the FRW setup must be stated explicitly:
\begin{enumerate}
\item \textbf{Conformally coupled massless scalar only.}
  Theorem~\ref{thm:FRW-typeII} applies to $\xi = 1/6$, $m = 0$.
  For $m > 0$ or $\xi \ne 1/6$, the field equation is not conformally
  covariant; the unitary $U$ does not intertwine the algebras; the
  type of the crossed product is unknown.
\item \textbf{Bounded-away-from-singularity assumption.}
  The condition $\eta_i > 0$ (diamond bounded away from the Big Bang)
  is essential. When $\eta_i = 0$, the conformal rescaling $f \mapsto a^3 f$
  degenerates and the isomorphism breaks.
\item \textbf{No quantitative $C = V$ dictionary.}
  Unlike the DSSYK/JT setting of~\cite{HPS2024}, there is no known
  gravity dual for the FRW diamond that would give a quantitative
  complexity-equals-volume formula. All computations below are at the
  level of the modular Lanczos recursion, without a holographic check.
\item \textbf{ECI v6.0.44 explicit caveat.}
  Section 8 of the structural limitations of ECI v6.0.44~\cite{ECI_v6}
  states: ``The type of $\mathcal{A}(\mathcal{O}) \rtimes_\sigma \mathbb{R}$
  for $\mathcal{O}$ a causal patch of a non-stationary FRW observer
  (radiation- or matter-dominated era) is, to our knowledge, not derived
  in the published math.OA literature.''
  The gap is filled by~\cite{FRWnote}, submitted 2026-05-02, but has
  not yet been refereed. We proceed on the basis of the technical
  arguments of~\cite{FRWnote} while acknowledging pre-refereeing status.
\end{enumerate}

%% ─────────────────────────────────────────────────────────────────────
\section{Krylov Complexity in the FRW Diamond}
\label{sec:krylov-FRW}
%% ─────────────────────────────────────────────────────────────────────

\subsection{Modular Krylov construction}

The Krylov complexity of an operator $\mathcal{O}$ under the modular flow
$\sigma_s^\FRW$ is defined via the modular Liouvillian
$\mathcal{L}_\FRW = [K_\FRW, \cdot]$.
The Krylov basis is generated by the Gram--Schmidt orthonormalisation
of $\{\mathcal{L}_\FRW^n \mathcal{O}\}_{n \ge 0}$ with respect to the
KMS inner product
\begin{equation}
  (\mathcal{O}_1, \mathcal{O}_2)_\beta
  \;:=\; \int_0^\beta \frac{d\tau}{\beta}
  \langle \mathcal{O}_1^\dagger\,
  e^{-(\beta-\tau) K_\FRW}\,\mathcal{O}_2\,
  e^{-\tau K_\FRW}\rangle_\FRW,
  \label{eq:KMS-inner-product}
\end{equation}
where $\langle \cdot \rangle_\FRW$ is the vacuum expectation value in
the Bunch--Davies state.
The Lanczos coefficients $\{b_n^\FRW\}$ are defined by the three-term
recurrence in this inner product.

\subsection{The autocorrelation function and Lanczos asymptotics}

The Lanczos coefficients are determined by the modular autocorrelation
function
\begin{equation}
  C^\FRW(s) \;:=\; (\mathcal{O},\,\sigma_s^\FRW(\mathcal{O}))_\beta
  \;=\; \langle \mathcal{O}^\dagger\,\sigma_s^\FRW(\mathcal{O})\rangle_\FRW.
  \label{eq:autocorr}
\end{equation}
By Theorem~\ref{thm:FRW-typeII}(1), this equals
\begin{equation}
  C^\FRW(s)
  \;=\; \langle U\mathcal{O}U^{-1},\,\sigma_s^\HL(U\mathcal{O}U^{-1})\rangle_\Mink
  \;=\; C^\HL(s),
  \label{eq:autocorr-pullback}
\end{equation}
where $C^\HL(s)$ is the Hislop--Longo autocorrelation function for
the operator $\tilde{\mathcal{O}} := U\mathcal{O}U^{-1}$ in the
Minkowski diamond vacuum.
In other words, \emph{the modular autocorrelation function in FRW is
identical to the Hislop--Longo autocorrelation function in Minkowski,
evaluated on the conformally related operator.}

For the free massless scalar, the Hislop--Longo two-point function
on the Minkowski diamond at $\beta = 2\pi$ decays as
\begin{equation}
  C^\HL(s) \;\sim\; \frac{1}{(2\sinh(\pi s))^{d-2}}
  \;\xrightarrow{d=4}\; \frac{1}{(2\sinh(\pi s))^2}
  \;\sim\; e^{-2\pi|s|}
  \quad (|s| \to \infty),
  \label{eq:autocorr-HL}
\end{equation}
giving $C^\FRW(s) = C^\HL(s)$ with the same exponential decay.
By the universal relationship between the spectral density of $C(s)$
and the Lanczos asymptotics (Parker et al.~\cite{Parker2019}),
an exponential decay $C(s) \sim e^{-2\pi|s|}$ corresponds to a
spectral density supported on $[-2\pi, 2\pi]$ (an arc), giving
linear Lanczos growth:
\begin{equation}
  b_n^\FRW \;\sim\; \frac{\pi}{\beta}\,n
  \;=\; \frac{\pi}{2\pi}\,n \;=\; \frac{n}{2}
  \quad (\beta = 2\pi).
  \label{eq:lanczos-FRW}
\end{equation}

\begin{remark}[Caution: free field vs. interacting]
\label{rem:free-field}
The estimate~\eqref{eq:autocorr-HL} and its consequence~\eqref{eq:lanczos-FRW}
are for the \emph{free} massless scalar, which is \emph{not} maximally
chaotic in the SYK sense.
In free field theories, the Lanczos coefficients typically grow as
$b_n \sim n^{1/2}$ (not linearly), leading to polynomial rather than
exponential Krylov complexity growth.
[IMPORTANT CAVEAT: the free-field autocorrelation decay rate does not
generically imply linear Lanczos growth; the relationship between
spectral density support and Lanczos asymptotics is more subtle.
Parker et al.~\cite{Parker2019} show linear Lanczos growth for
interacting and chaotic systems; for free fields the situation is
different. The estimate~\eqref{eq:lanczos-FRW} should be regarded
as a formal analogy, not a rigorous derivation.]
The step from $C^\FRW(s) = C^\HL(s)$ to $\lK^\FRW = 2\pi/\beta$
requires either: (a) passing to an interacting field in the FRW
diamond (breaking exact conformal invariance), or (b) invoking the
operator growth hypothesis as applying to the modular flow on the
type-II$_\infty$ algebra in a ``chaotic'' sense to be made precise.
This is the principal open gap.
\end{remark}

\subsection{Proposed Krylov complexity growth rate}

Accepting the formal analogy and the modular Krylov construction of
Caputa et al.~\cite{Caputa2024}, the Krylov complexity of an operator
$\mathcal{O}$ under the FRW modular flow grows as
\begin{equation}
  \CK^\FRW(s) \;\sim\; e^{\lK^\FRW\,|s|},
  \qquad |s| \to \infty,
  \label{eq:Krylov-FRW-growth}
\end{equation}
with the Krylov Lyapunov exponent
\begin{equation}
  \lK^\FRW \;=\; 2b_1^\FRW \;\approx\; \frac{2\pi}{\beta}
  \;=\; \frac{2\pi}{2\pi} \;=\; 1,
  \qquad (\beta = 2\pi).
  \label{eq:Krylov-FRW-rate}
\end{equation}
This is the same rate as in the Rindler wedge~\cite{ModularShadow}
and in the SYK model~\cite{Parker2019, HPS2024}, and saturates the
MSS bound~\eqref{eq:mss} exactly.

%% ─────────────────────────────────────────────────────────────────────
\section{The Conjecture and Supporting Arguments}
\label{sec:conjecture}
%% ─────────────────────────────────────────────────────────────────────

\begin{conjecture}[FRW Krylov Saturation]
\label{conj:FRW-krylov}
\emph{[CONJECTURE]}
Let the setup be as in Theorem~\ref{thm:FRW-typeII}: conformally
coupled massless scalar on radiation- or matter-dominated FRW, in the
conformal vacuum, on a causal diamond bounded away from the Big Bang
($\eta_i > 0$).
Then the Krylov complexity of operators evolved by the FRW modular
flow $\sigma_s^\FRW$ grows exponentially at late modular times with
Lyapunov exponent
\begin{equation}
  \lK^\FRW \;=\; \frac{2\pi}{\beta_\FRW} \;=\; 1,
  \label{eq:conj}
\end{equation}
where $\beta_\FRW = 2\pi$ is the KMS inverse temperature of the
Bunch--Davies vacuum with respect to the FRW modular flow.
This saturates the MSS chaos bound~\eqref{eq:mss} exactly.
\end{conjecture}

\subsection{Supporting arguments}

\paragraph{A1. Algebraic universality.}
By Theorem~\ref{thm:FRW-typeII}, the FRW diamond algebra is type II$_\infty$
with Connes spectrum $\Gamma = \mathbb{R}$ and modular flow isomorphic
(via $U$) to the Hislop--Longo flow.
The companion paper~\cite{ModularShadow} conjectures that MSS saturation
is a structural property of type-II$_\infty$ algebras at KMS temperature
$\beta = 2\pi$: the Bisognano--Wichmann surface gravity $\kappa_R = 2\pi T_H$
appears as a structural constant fixed by the algebraic data.
If this conjecture is correct, Conjecture~\ref{conj:FRW-krylov} follows
as an instance in the FRW setting.

\paragraph{A2. Autocorrelation function argument.}
By~\eqref{eq:autocorr-pullback}, $C^\FRW(s) = C^\HL(s)$.
The Hislop--Longo autocorrelation decays as $e^{-2\pi|s|}$ in $d=4$
(at $\beta = 2\pi$), which, if one applies the Parker et al.\ universal
operator growth hypothesis~\cite{Parker2019} to the modular flow,
would give $b_n \sim n/2$ and $\lK = 2\pi/\beta = 1$.
The main gap (Remark~\ref{rem:free-field}) is that this argument
requires the interacting/chaotic version of the growth hypothesis
to apply to modular dynamics.

\paragraph{A3. Consistency with the DSSYK holographic result.}
Heller, Papalini, and Schuhmann~\cite{HPS2024} find $\lK = 2\pi/\beta$
in DSSYK at finite temperature.
The DSSYK model at large $N$ (in the triple-scaled limit) approaches
a type-II$_\infty$ structure, and the KMS temperature of the Hartle--Hawking
state in the gravity dual is identified with the SYK inverse temperature $\beta$.
If Conjecture~\ref{conj:FRW-krylov} is correct, it provides an independent
(non-holographic) derivation of the same rate in a different type-II$_\infty$
algebra --- the FRW diamond --- giving universal evidence that
$\lK = 2\pi/\beta$ is an algebraic invariant, not a holographic artefact.

\paragraph{A4. Modular shadow identification.}
From~\eqref{eq:eci-mss} and the Faulkner--Speranza identification
$\sgen \leftrightarrow \Smod$ (flagged as a conjecture), the ECI
saturation rate $2\pi T_H$ equals $2\pi/\beta$ at $T_H = 1/(2\pi)$
(in units $k_B = \hbar = 1$, $\beta = 2\pi$).
This is exactly the rate~\eqref{eq:conj} in Conjecture~\ref{conj:FRW-krylov}.
The FRW diamond realization provides a new geometrically distinct instance
of this universal rate, strengthening the evidence for the modular shadow
identification.

\subsection{The principal obstruction: free field vs. chaos}

The main obstruction to proving Conjecture~\ref{conj:FRW-krylov} is the
free-field nature of the setup.
For free fields in Minkowski space, the modular Krylov complexity under
the Bisognano--Wichmann boost flow does \emph{not} generically exhibit
exponential growth with $\lK = 2\pi/\beta$; the free-field Lanczos
coefficients grow sub-linearly and Krylov complexity grows polynomially
or as a power law~\cite{Parker2019}.
The exponential Krylov growth at $\lK = 2\pi/\beta$ is a property
of \emph{maximally chaotic} interacting systems (SYK, near-extremal
black holes), not of free fields.

The Araki relative modular operator approach to the $\sgen \leftrightarrow \Smod$
bridge in type-II$_\infty$ algebras is pursued in~\cite{EREPR_Araki}.

Three routes to resolve this:
\begin{enumerate}
\item[(R1)] \textbf{Interacting FRW fields.}
  Introduce a $\lambda\phi^4$ interaction on the FRW background.
  The conformally coupled scalar with $\lambda\phi^4$ is not conformally
  invariant in 4d (the $\beta$-function is non-trivial).
  However, the \emph{type} of the crossed product (type II$_\infty$)
  is expected to persist under weak interactions (it is a topological
  property of the modular flow), and the Lanczos growth rate might
  approach $\pi/\beta$ in the strongly interacting (renormalized) theory.
  This requires an FRW analogue of the SYK large-$N$ argument.

\item[(R2)] \textbf{Large-$N$ limit in FRW.}
  Consider $N$ copies of the conformally coupled scalar on the FRW
  diamond, with all-to-all $\phi^4$-type couplings (SYK-like random
  coupling).
  In the large-$N$ limit, one expects (by analogy with SYK~\cite{SYK_Kitaev})
  a type-II$_\infty$ modular flow with $\lK \to 2\pi/\beta$.
  This is the most direct route and the most computationally tractable.

\item[(R3)] \textbf{Algebraic definition of ``chaotic'' modular flow.}
  Define a notion of ``maximally chaotic modular automorphism'' in
  a type-II$_\infty$ algebra (e.g., via the Connes entropy or relative
  modular growth), and show that the Hislop--Longo flow, when restricted
  to an appropriate sub-algebra, falls in this class.
  This is the most mathematically ambitious route and the natural
  direction for the companion modular shadow paper~\cite{ModularShadow}.
\end{enumerate}

%% ─────────────────────────────────────────────────────────────────────
\section{Open Questions and Research Programme}
\label{sec:open}
%% ─────────────────────────────────────────────────────────────────────

\begin{question}\label{q:free-vs-interacting}
  Does the modular Krylov complexity of the \emph{free} conformally
  coupled scalar on the FRW diamond grow exponentially or polynomially?
  If polynomially, at what power?
  (This is an explicit computation using~\cite{Parker2019}
  applied to $C^\HL(s) \sim (2\sinh(\pi s))^{-2}$.)
\end{question}

\begin{question}\label{q:large-N}
  Do $N$ conformally coupled scalars on the FRW diamond with large-$N$
  SYK-type random couplings satisfy $\lK \to 2\pi/\beta$ as $N \to \infty$?
  What is the $O(1/N)$ correction to $\lK$?
  (This mirrors the SYK analysis of~\cite{Parker2019} and the quantum
  correction of~\cite{HPS2024} in the holographic setting.)
\end{question}

\begin{question}\label{q:dS}
  Does the same conjecture hold for the de Sitter static patch?
  In the static patch, the KMS temperature is $\beta_\dS = 2\pi/H$
  (Gibbons--Hawking), giving $\lK^\dS = H$ (the Hubble rate).
  The modular shadow conjecture then predicts $\lK = H$, matching the
  de Sitter ``scrambling rate'' and connecting to cosmic censorship.
\end{question}

\begin{question}\label{q:holo-FRW}
  Is there a two-dimensional gravity dual for the FRW diamond that would
  give a $C = V$ formula analogous to HPS~\cite{HPS2024}?
  The ``FRW/CFT'' correspondence (if it exists) would provide the
  quantitative check. A candidate is the celestial holography on the
  FRW past light cone, but no quantitative $C = V$ formula is known.
\end{question}

\begin{question}\label{q:bigbang}
  What happens to $\lK^\FRW$ as $\eta_i \to 0$?
  The Big Bang singularity is expected to introduce new structure
  (possibly a transition from type II$_\infty$ to type II$_1$ or
  type III, mirroring the singularity theorem in the modular-flow setting).
  If $\lK$ diverges as $\eta_i \to 0$, this could be a quantum-gravity
  signature of the Big Bang.
\end{question}

\begin{question}\label{q:observable}
  Is $\lK^\FRW$ observable in the FRW setting?
  A possible route: analog gravity experiments simulating conformally
  coupled fields in a FRW-like background (e.g., acoustic Hawking radiation
  in an expanding BEC, cf.\ Steinhauer 2019 cited in~\cite{ECI_v6}).
  The ECI Bisognano--Wichmann window $u_\mathrm{max} = 2\pi$ constrains
  the temporal window in which modular chaos is observable.
\end{question}

\paragraph{Proposed computational programme ($\sim 4$ weeks).}
\begin{enumerate}
\item \textit{Week 1.} Explicit computation of $C^\FRW(s)$ for the
  free massless scalar via~\eqref{eq:autocorr-pullback} and the known
  Hislop--Longo two-point function; Lanczos coefficients $b_n^\FRW$
  for $n = 1, \ldots, 50$ via numerical recursion.
  Test whether $b_n^\FRW$ grows linearly ($\lK = 1$) or sub-linearly.

\item \textit{Week 2.} $N = 2, 4, 8$ conformally coupled scalars on
  a finite FRW diamond with SYK-type random $\phi^4$ couplings.
  Compute OTOC Lyapunov exponent and check $\lL \le 2\pi/\beta$.
  Does $\lL \to 2\pi/\beta$ as $N$ or coupling increases?

\item \textit{Week 3.} Analytical approximation of $b_n^\FRW$ in the
  large-$n$ limit using the spectral density of $K_\FRW$ via the
  Kac--Moody generator conjugation~\eqref{eq:K-FRW}.
  Attempt asymptotic analysis using the known M\"{o}bius-flow spectral
  measure on the diamond.

\item \textit{Week 4.} Write-up; connection to the de Sitter case
  (Question~\ref{q:dS}); comparison with DSSYK quantum corrections
  of~\cite{HPS2024}; submission as Comm.\ Math.\ Phys.\ paper
  (or companion preprint to~\cite{ModularShadow}).
\end{enumerate}

%% ─────────────────────────────────────────────────────────────────────
\section*{Acknowledgments}
%% ─────────────────────────────────────────────────────────────────────

This note was drafted as Opus Bridge \#16, extending the
DSSYK/Krylov holographic complexity analysis of~\cite{HPS2024}
to the FRW diamond, using the type-II$_\infty$ classification
of~\cite{FRWnote} and the modular shadow framework of~\cite{ECI_v6,ModularShadow}.
The [CONJECTURE] tags throughout reflect the current status of the
central claim (Conjecture~\ref{conj:FRW-krylov}): supported by structural
arguments but not proven.

%% ─────────────────────────────────────────────────────────────────────
\begin{thebibliography}{99}
%% ─────────────────────────────────────────────────────────────────────

\bibitem{HPS2024}
M.~P.~Heller, J.~Papalini, and T.~Schuhmann,
``Krylov spread complexity as holographic complexity beyond JT gravity,''
\textit{Phys. Rev. Lett.} \textbf{135} (2025) 151602,
\href{https://doi.org/10.1103/PhysRevLett.135.151602}{doi:10.1103/PhysRevLett.135.151602},
\href{https://arxiv.org/abs/2412.17785}{arXiv:2412.17785 [hep-th]}.
[Verbatim abstract verified via WebFetch on 2026-05-04.]

\bibitem{Parker2019}
D.~E.~Parker, X.~Cao, A.~Avdoshkin, T.~Scaffidi, and E.~Altman,
``A Universal Operator Growth Hypothesis,''
\textit{Phys. Rev. X} \textbf{9} (2019) 041017,
\href{https://doi.org/10.1103/PhysRevX.9.041017}{doi:10.1103/PhysRevX.9.041017},
\href{https://arxiv.org/abs/1812.08657}{arXiv:1812.08657 [cond-mat.str-el]}.
[Verbatim abstract verified via WebFetch on 2026-05-04.
Key result: $\lambda_L \le 2\alpha$ where $b_n \sim \alpha n$;
$\alpha = \pi/\beta$ at MSS saturation in SYK numerics.]

\bibitem{MSS2016}
J.~Maldacena, S.~H.~Shenker, and D.~Stanford,
``A bound on chaos,''
\textit{JHEP} \textbf{08} (2016) 106,
\href{https://doi.org/10.1007/JHEP08(2016)106}{doi:10.1007/JHEP08(2016)106},
\href{https://arxiv.org/abs/1503.01409}{arXiv:1503.01409 [hep-th]}.

\bibitem{FRWnote}
K.~Remondi\`{e}re,
``Type II$_\infty$ for the gravitational algebra of a comoving observer
in radiation-dominated FRW via conformal pullback,''
Technical note, 2026-05-02,
\url{https://github.com/AIdevsmartdata/crossed-cosmos/paper/frw_typeII_note/frw_note.tex}.
[Main theorems: Theorem~1 (radiation-dominated FRW), Theorem~2
(matter-dominated FRW), Corollary (slow-roll de Sitter).
Sympy-verified: conformal-covariance identity and symplectic-form cancellation.
Status: pre-refereeing.]

\bibitem{FRWconformal}
ECI sub-agent (Sonnet 4.6),
``Conformal Pullback of the Modular Flow in Radiation-Dominated FRW:
Explicit Formula and Spectrum Analysis,''
Technical note, 2026-05-02,
\url{https://github.com/AIdevsmartdata/crossed-cosmos/notes/frw_conformal_pullback_attack_2026_05_02.md}.

\bibitem{ModularShadow}
[ECI author(s)],
``The Modular Shadow Conjecture:
MSS Chaos-Bound Saturation as a Signature of Type-II$_\infty$
Crossed-Product Structure,''
Submitted to \textit{Letters in Mathematical Physics}, 2026-05-04.
\url{https://github.com/AIdevsmartdata/crossed-cosmos/notes/eci_v7_aspiration/MODULAR_SHADOW/modular_shadow_LMP.tex}.

\bibitem{ECI_v6}
[ECI author(s)],
``Entanglement--Complexity--Information: a phenomenological architecture
for quantum gravity phenomenology (v6.0.44),''
Zenodo record 20021358, 2026-05-04.
[Cited for \S sec:mss-shadow (modular shadow conjecture) and
\S structural limitations item 8 (FRW non-stationary classification gap).]

\bibitem{CLPW2023}
V.~Chandrasekaran, R.~Longo, G.~Penington, and E.~Witten,
``An algebra of observables for de~Sitter space,''
\textit{JHEP} \textbf{02} (2023) 082,
\href{https://doi.org/10.1007/JHEP02(2023)082}{doi:10.1007/JHEP02(2023)082},
\href{https://arxiv.org/abs/2206.10780}{arXiv:2206.10780 [hep-th]}.

\bibitem{Witten2022}
E.~Witten,
``Gravity and the crossed product,''
\textit{JHEP} \textbf{10} (2022) 008,
\href{https://doi.org/10.1007/JHEP10(2022)008}{doi:10.1007/JHEP10(2022)008},
\href{https://arxiv.org/abs/2112.12828}{arXiv:2112.12828 [hep-th]}.

\bibitem{FaulknerSperanza2024}
T.~Faulkner and A.~J.~Speranza,
``Gravitational algebras and the generalized second law,''
\textit{JHEP} \textbf{11} (2024) 099,
\href{https://doi.org/10.1007/JHEP11(2024)099}{doi:10.1007/JHEP11(2024)099},
\href{https://arxiv.org/abs/2405.00847}{arXiv:2405.00847 [hep-th]}.

\bibitem{Caputa2024}
P.~Caputa, J.~M.~Mag\'{a}n, D.~Patramanis, and E.~Tonni,
``Krylov complexity of modular Hamiltonian evolution,''
\textit{Phys. Rev. D} \textbf{109} (2024) 086004,
\href{https://doi.org/10.1103/PhysRevD.109.086004}{doi:10.1103/PhysRevD.109.086004},
\href{https://arxiv.org/abs/2306.14732}{arXiv:2306.14732 [hep-th]}.

\bibitem{HislopLongo1982}
P.~D.~Hislop and R.~Longo,
``Modular structure of the local algebras associated with the free
massless scalar field theory,''
\textit{Commun. Math. Phys.} \textbf{84} (1982) 71--85.

\bibitem{BisognanoWichmann1976}
J.~J.~Bisognano and E.~H.~Wichmann,
``On the duality condition for a Hermitian scalar field,''
\textit{J. Math. Phys.} \textbf{17} (1976) 303.

\bibitem{SYK_Kitaev}
A.~Kitaev, ``A simple model of quantum holography,''
Talks at KITP, 7 and 23 April 2015,
\url{http://online.kitp.ucsb.edu/online/entangled15/kitaev/}.

\bibitem{BIX_lemma_A1}
K.~Remondi\`{e}re,
``Appendix A: Closure of the Eigenvalue-Crossing Gap for the Hadamard SLE on Bianchi~IX,''
arXiv:[ECI v6.2 Bianchi~IX, in preparation, 2026].
%% [Fix \#9: BIX cite added to FRW-setup for Hadamard/type-II analogy.]

\bibitem{EREPR_Araki}
[ECI author(s)],
``Araki Relative Modular Operator in Type-II$_\infty$: ER=EPR Reopen (A13-2),''
arXiv:[ECI v6.2 ER=EPR route, in preparation, 2026].
%% [Fix \#10: added as companion paper pursuing S_gen <-> C_Krylov bridge via Araki.]

\end{thebibliography}

\end{document}
